\documentclass[11pt]{article}

\usepackage[letterpaper, margin=1in]{geometry}
\usepackage[default]{opensans}
\usepackage[skip=5pt]{parskip}
\usepackage{setspace}
\usepackage{enumitem}

\title{BRICK Bylaws}

\newcounter{h1}
\newcounter{h2}[h1]
\newcounter{h3}[h2]

\begin{document}
\newcommand{\tab}[1]{\stepcounter{h1}\par\large\textbf{\arabic{h1}. #1}\normalsize\par}
\newcommand{\subtab}[1]{\stepcounter{h2}\par\textbf{\arabic{h1}.\arabic{h2}. #1}\par}
\newcommand{\subsubtab}[1]{\stepcounter{h3}\par\small\textbf{\arabic{h1}.\arabic{h2}.\arabic{h3}. #1}\normalsize\par}
\renewcommand{\labelenumi}{\arabic{enumi}.} 
\renewcommand{\labelenumii}{\arabic{enumi}.\arabic{enumii}.}
\renewcommand{\labelenumiii}{\arabic{enumi}.\arabic{enumii}.\arabic{enumiii}.}
\renewcommand{\labelenumiv}{\arabic{enumi}.\arabic{enumii}.\arabic{enumiii}.\arabic{enumiv}.}

\setstretch{1.5}
\setenumerate{noitemsep}
\setitemize{noitemsep}
\raggedright


\begin{center}
\Large
\textbf{BRICK Bylaws}

\normalsize
Revision 1.2 (February 25, 2025)
\end{center}


\tab{Name and purpose}

\subtab{Name}

This organization shall be known as BRICK or Brick.

\subtab{Mission statement}

BRICK is a student organization with the objective of mobilizing CMU's community toward activism centered on queer issues, by providing students with support, resources, and community. Our goal is to secure the rights and safety of queer individuals and oppose oppression in all forms.

\subtab{Purpose}

\begin{enumerate}
  \item The purpose of BRICK shall be:  
  \begin{enumerate}
    \item To organize, support, and provide resources on protests or other activities which support any issue of importance to CMU's queer community.  
    \item To provide support and mutual aid within CMU's queer community and to its allies.  
    \item To foster a campus community inclusive of queer and marginalized groups.  
  \end{enumerate}
  \item BRICK shall not operate for profit and no income of the organization may be distributed to private individuals.  
  \item All activities and functions of the organization must be legal under University, local, state, and federal laws.
\end{enumerate}

\tab{Definitions}

\begin{enumerate}
  \item \textbf{Acclamation:} A vote by acclamation shall be defined as a procedure which decides a question by unanimous consent of the attendees of a COM.  
  \item \textbf{Ballot item:} A ballot item shall be defined as a question to be decided by a formal vote.  
  \item \textbf{Collective organization:} The collective organization shall be defined as the organization's general membership, in contrast to any of its specific committees or interest groups.  
  \item \textbf{Collective organization meeting:} A collective organization meeting, or COM, shall be defined as a meeting which all members may attend as a right of membership. The primary purpose of COMs is to make decisions requiring the consensus of the organization at large. This definition is duplicated in section 5.2 ("collective organization meetings").  
  \item \textbf{Committee:} A committee shall be defined as a group which coordinates activities relating to a specific, narrow, and time-sensitive purpose. This definition is duplicated in section 6.1 ("structure of committees").  
  \item \textbf{Committee meeting:} A committee meeting shall be defined as a meeting at which one-half or more of a committee's chairs are present and which all at-large members of the committee may attend as a right of membership. This definition is duplicated in section 6.4 ("committee meetings").
  \item \textbf{Day:} Unless specified, a day shall be defined as a calendar day, such that "within 1 day" of a designated point in time shall include any time during the same or the following calendar day.  
  \item \textbf{Election:} An election shall be defined as a procedure consisting of a nomination period and a formal vote, which shall be used to select officers.  
  \begin{enumerate}
    \item \textbf{General election:} A general election shall be defined as a semesterly election which shall fill all officer positions.  
    \item \textbf{Special election:} A special election shall be defined as an irregularly scheduled election which shall fill one or more vacant officer positions.  
  \end{enumerate}
  \item \textbf{Formal vote:} A formal vote shall be defined as an anonymous process to decide one or more questions through an in-person election. Formal votes shall occur during COMs. Most COMs will occur to resolve proposed resolutions from the previous COM, or for general or special elections of officers. This definition is duplicated in section 8.2.1 ("purpose of formal votes").  
  \item \textbf{Interest group:} An interest group, or IG, shall be defined as a group of individuals and one or more committees with a shared focus, with the purpose of sharing information and updates relating to that focus, rather than carrying out actions, which is the duty of committees. This definition is duplicated in section 7.1 ("interest groups").  
  \item \textbf{Majority:} A majority shall be defined as greater than one-half, discounting abstentions.  
  \item \textbf{Member:} A member shall be defined as an individual associated with the organization who fulfills all requirements of membership and receives all rights associated with such status.  
  \begin{enumerate}
    \item \textbf{Non-voting member:} A non-voting member shall be defined as a member who is not a voting member.  
    \item \textbf{Voting member:} A voting member shall be defined as a member with the right to vote in COMs and formal votes.  
  \end{enumerate}
  \item \textbf{Officer:} An officer shall be defined as any individual holding one or more officer positions.  
  \item \textbf{Officer board:} The officer board shall be defined as the group composed of all officers of the organization.  
  \item \textbf{Officer position:} An officer position shall be defined as an elected role held by one or more individuals which carries specific associated rights and duties.  
  \item \textbf{Poll:} A poll shall be defined as a ballot item in a formal vote which does not create a resolution, but may be used as a condition of another resolution or to seek the input of the collective organization on some matter.  
  \item \textbf{Prohibition from membership:} Prohibition from membership shall refer to a status, applying to individuals, prohibiting membership in the organization for some specified period of time, which may be indefinite. This term is synonymous with "suspension." This definition is duplicated in section 3.3 ("suspension and prohibition of membership").  
  \item \textbf{Resolution:} A resolution shall be defined as a formally passed decision of the collective organization, which mandates the organization to take a specific action or set of actions.  
  \begin{enumerate}
    \item \textbf{Simple resolution:} A simple resolution or majority resolution shall be defined as a resolution which may be passed by acclamation or by a majority formal vote ballot item.  
    \item \textbf{Supermajority resolution:} A supermajority resolution shall be defined as a resolution which may be passed only by a supermajority formal vote ballot item.  
  \end{enumerate}
  \item \textbf{Single Transferable Vote (STV)}: Single Transferable Vote, or STV, shall be defined as a vote counting method for ranked choice ballot items for multiple seats, which shall be counted using Meek's method.  
  \item \textbf{Special interest group:} A special interest group, or SIG, shall be defined as an interest group with a limited lifespan. This definition is duplicated in section 7.2 ("special interest groups").  
  \item \textbf{Supermajority:} A supermajority shall be defined as greater than two-thirds, discounting abstentions.  
  \item \textbf{Suspension:} Suspension shall refer to a status, applying to individuals, prohibiting membership in the organization for some specified period of time, which may be indefinite. This term is synonymous with "prohibition from membership." This definition is duplicated in section 3.3 ("suspension and prohibition of membership").  
  \item \textbf{Weak majority:} A weak majority shall be defined as greater than or equal to one-half, discounting abstentions.  
  \item \textbf{Written vote:} A written vote shall be defined as a procedure to decide a question among a group of fixed size, such as a committee or officer board. This definition is duplicated in section 8.1 ("written votes").
\end{enumerate}

\tab{Membership}

\subtab{Membership requirements}

\begin{enumerate}
  \item To be eligible for membership, an individual must currently be a student, staff, or faculty member of CMU.  
  \begin{enumerate}
    \item The collective organization may override this requirement conditionally or for specific individuals, possibly subject to some specified time constraint. No resolution may override ineligibility due to suspension from the organization.  
    \item Member status is lost immediately upon loss of eligibility.  
  \end{enumerate}
  \item In order for an individual eligible for membership to become a non-voting member, they must follow an onboarding procedure established by the collective organization.  
  \begin{enumerate}
    \item BRICK is committed to abiding by the CMU Statement of Assurance.  
    \item Hazing will not be used as a condition of membership in this organization.  
    \item The onboarding process may not require dues or any form of financial commitment.  
    \item Any individual who is eligible for membership must be able to complete onboarding and to do so within a reasonable time period. An individual may only be rejected from the onboarding process through suspension from the organization.  
  \end{enumerate}
  \item In order for a non-voting member to become a voting member, they must (i) become a chair of a committee or (ii) become an officer.  
  \item If a voting member is (i) not a chair of any committees and (ii) is not an officer, for the full duration of one fall or spring semester, on the last day of that semester the voting member shall lose voting membership.  
  \item Members may withdraw their membership at any time.
\end{enumerate}

\subtab{Rights of members}

\begin{enumerate}
  \item All voting and non-voting members shall have the right to attend and fully participate in COMs, which shall include the right to introduce proposed resolutions.  
  \begin{enumerate}
    \item Only voting members shall have the right to participate in votes by acclamation.  
  \end{enumerate}
  \item All voting and non-voting members shall have the right to join any and all interest groups.  
  \item All voting and non-voting members shall have the right to committee chair candidacy.  
  \item All voting and non-voting members shall have the right to officer candidacy.  
  \item All voting members shall have the right to cast a ballot in all formal votes.
\end{enumerate}

\subtab{Suspension and prohibition of membership}

\begin{enumerate}
  \item The terms "suspended" or "prohibited from membership" shall identically refer to a status, applying to individuals, prohibiting membership in the organization for some specified period of time, which may be indefinite. This definition is duplicated in section 2 ("definitions").  
  \begin{enumerate}
    \item No part of this document shall use the term "suspended" or any variants thereof to refer to suspension from Carnegie Mellon University or any other bodies without appropriate clarification.  
  \end{enumerate}
  \item Membership may be suspended, or a suspension may be altered or undone, by a written majority vote of the officer board.  
  \begin{enumerate}
    \item No prior membership shall be required.  
    \item If the member to be suspended is an officer, removal shall occur by a written supermajority vote of the other officers. The officer is then immediately removed from all officer positions, which become vacant.  
  \end{enumerate}
\end{enumerate}

\tab{Officers}

\subtab{Purpose of officers}

The purpose of officers is to:

\begin{enumerate}
  \item Fulfill any duties imposed on the organization or its officers by Carnegie Mellon University, its Student Government, or other outside parties.  
  \item Neutrally facilitate the conditions under which the organization may make decisions regarding the use of its collective resources.  
  \item Function as contact points and information clearinghouses for individuals within or outside of the organization.  
  \item Oversee processes related to membership, which is to include handling issues of a confidential nature relating to members and membership status.
\end{enumerate}

\subtab{Officer positions}

\subsubtab{Clerks}

There may be multiple clerks at any given time. Clerks are responsible for chairing and overseeing collective organization meetings, which shall include the following responsibilities:

\begin{itemize}
  \item Conducting COM in accordance with all bylaws and all applicable procedures instituted through resolutions of the collective organization.
  \item Conducting business neutrally and giving all attendees an equitable chance to be heard.
  \item Ensuring minutes are taken and that the minutes-taker(s) record discussion in a neutral manner.
  \item Overseeing the reservation of rooms or locations for the organization's use.
\end{itemize}

\subsubtab{Finance contact point}

The finance contact point is responsible for overseeing:

\begin{itemize}
  \item The budget and finances of the organization, including any financial responsibilities imposed by CMU or Student Government.
  \item Space allocation and access for the organization.
  \item Storage of the collective organization's physical possessions.
\end{itemize}

The finance contact point is not responsible for approving or denying spending by the organization unless delegated such powers.

\subsubtab{Communications contact point}

The communications contact point's responsibilities consist of:

\begin{enumerate}
  \item Managing access to the organization's public digital presences, including on TartanConnect and social media platforms.  
  \item Managing access to mailing lists or other communications issued by the collective organization.  
  \item Functioning as a contact point regarding use of the organization's name and branding.
\end{enumerate}

\subsubtab{Campus contact point}

The campus contact point is responsible for representing BRICK within CMU, particularly in cooperation with other student organizations. The campus contact point should serve as required as a resource for recruiting and onboarding activities.

\subsubtab{Documents contact point}

The documents contact point is responsible for:

\begin{enumerate}
  \item Storing physical and digital documents produced within the organization, including agendas, notes, and statements, and making such documents available for access as required by the organization.  
  \item Overseeing information security within the organization, including providing security recommendations or regulations for internal communications, organizing routine security compliance checks, and coordinating responses to unintentional disclosure of confidential information.  
  \item Overseeing archival efforts, including transferral of documents to the University Archives should they become suitable for public disclosure.
\end{enumerate}

\subsubtab{Local community contact point}

The local community contact point is responsible for:

\begin{enumerate}
  \item Maintaining an awareness of off-campus issues, protests, and other events, particularly within Pittsburgh.  
  \item Maintaining relationships with off-campus organizations.
\end{enumerate}

\subtab{Selection of officers}

\subsubtab{General and special elections}

Two types of election shall occur: % Note: I accidentally deleted the (somewhat paranoid/implied) provision that these are the *only* means by which an officer position may be granted. Amend later :-)

\begin{enumerate}
  \item General elections, by which the full officer board is re-selected. Precisely two general elections must occur per year:
  \begin{enumerate}
    \item One must occur no earlier than the first Friday after fall break and no later than the second-to-last Friday of the fall semester.
    \item One must occur no earlier than the first Friday after spring break and no later than the second-to-last Friday of the spring semester.
  \end{enumerate}
  \item Special elections, by which all vacant officer positions are filled. Special elections must be scheduled as required by section 4.4 ("vacancies").
\end{enumerate}

\subsubtab{Nominations}

Nominations for officer positions in a general or special election shall occur by the following procedure:

\begin{enumerate}
  \item A voting board should be assembled following the rules specified in section 8.2 ("formal votes").  
  \item A minimum of 14 days from the planned formal vote, the voting board must announce the times and locations of all COMs which will include time for nominations, and of the formal vote, as precisely as such information can be reasonably known.  
  \item At all COMs which include time for nominations, of which at least 2 must occur prior to the formal vote, the following procedure shall be followed:  
  \begin{enumerate}
    \item The current officers and/or voting board must provide descriptions of the roles of all officer positions which are up for election.  
    \item The voting board must list all current nominees. Any absentee nominations should be seconded during this time.  
    \item Nominations may be submitted or withdrawn for any officer position which is up for election.  
    \begin{enumerate}
      \item In-person nominations must be seconded immediately.  
    \end{enumerate}
  \end{enumerate}
  \item Only self-nominations shall be accepted.  
  \item Nominees must meet one of the following eligibility requirements:  
  \begin{enumerate}
    \item The nominee is a member of the organization.  
    \item One current officer or two current voting members are willing to vouch for the nominee, subject to a veto by a simple resolution or a written weak majority of the officer board.  
  \end{enumerate}
  \item Absentee nominations shall be open during the period following the first COM during which nominations are taken and prior to the last COM during which nominations are taken.  
  \item During the last COM prior to the formal vote, any candidate shall have the right to give up to one speech per position they are nominated for, and any voting member shall have the right to direct questions at any candidate. The voting board may set reasonable time limits for speeches, questions, or answers.  
  \begin{enumerate}
    \item At the discretion of the voting board, speeches and questions may take place during the nomination period in an earlier COM.  
    \item Absentee candidates may submit statements or answers to questions, subject to the discretion of the voting board.
  \end{enumerate}
\end{enumerate}

\subsubtab{Formal vote}

\begin{enumerate}
  \item The election shall occur through a formal vote.  
  \item If one or more clerks are to be elected, the formal vote shall contain one approval voting ballot item for clerks.  
  \begin{enumerate}
    \item The ballot item shall include a list of candidates, each with a checkbox or other unambiguous approve/disapprove marker.  
  \end{enumerate}
  \item The formal vote shall contain one ranked choice ballot item for each type of officer position which is up for election, with the exception of clerk.  
  \begin{enumerate}
    \item The ballot item shall include a list of candidates with a space in which a number may be written.  
    \item Voters should rank zero or more candidates using incrementing integers, starting with 1, which shall correspond to the most desirable candidate.  
  \end{enumerate}
  \item For officer positions other than clerk, vote counting shall occur by Instant Runoff Voting.  
  \begin{enumerate}
    \item Any ballots for which no choices are ranked, or for which all ranked choices are eliminated, shall be considered a vote for no candidate.  
    \item A candidate must receive a majority of votes to be elected. If this does not occur, the position shall be left vacant.  
  \end{enumerate}
  \item Officer positions, with the exception of clerks, shall be decided in the order they are listed in section 4.2 ("officer positions"). No individual may hold more than one non-clerk officer position.  
  \item Clerks shall be decided through approval voting. Any candidates with less than majority approval shall be eliminated, and highest approval shall be used to select the top 3 candidates among those remaining, or as many as is possible, whichever is fewer.  
  \item No individual may hold more than one clerk position simultaneously.
\end{enumerate}

\subtab{Transfer of positions}

\begin{enumerate}
  \item Following a fall semester general election, transfer of officer positions shall typically occur on the first day of the spring semester.
  \item Following a spring semester general election, transfer of officer positions shall typically occur on the last day of the spring semester.
  \item Following a special election, transfer of previously vacant officer positions shall typically occur immediately after vote counting is finished.
  \item Any occupied officer position not filled during an election shall become vacant upon the transfer of positions.
  \item Only with the consent of the recipient of an officer position (if any) and the full current officer board may the transfer of the position deviate from the typical scheduling. During a delayed transfer of one or more officer positions, the individual who formerly held a position may hold additional officer positions simultaneously even where prohibited by other clauses of these bylaws. However, the individual must remain a member of the organization.
\end{enumerate}

\subtab{Vacancies}

If an officer position is vacant, the duties of the position should be distributed among other officers. A special election for any vacant positions must be scheduled as soon as is reasonably possible, unless a general election is or will be scheduled within 28 days.

\subtab{Removal}

An officer may be removed from one or more of their officer positions by any of the following processes:

\begin{enumerate}
  \item The other officers may present a written supermajority vote.  
  \item A supermajority resolution may be passed by the collective organization with a minimum notice period of 72 hours.
\end{enumerate}

\tab{Collective organization}

\subtab{Rights and duties of the collective organization}

\begin{enumerate}
  \item Most activities of the organization shall be conducted through committees. The collective organization shall meet to create and oversee committees, to oversee use of resources belonging to the collective organization, and to regulate internal communications and security.  
  \item The collective organization shall reserve the right to carry out any action which is specified to occur through a resolution by any part of this document. Rights given to the collective organization throughout this document shall not be assumed to be delegable.  
  \item The collective organization shall have the right to pass simple resolutions to set the rules governing COMs or to set the time or location of COMs. This is a delegable right.  
  \item The collective organization shall have the right to pass simple resolutions to define the onboarding procedure applicable to new members. This is a delegable right.  
  \item The collective organization shall have the right to pass simple resolutions to oversee the scheduling of formal votes, to oversee the composition of voting boards, and to call for poll ballot items. This is a delegable right.  
  \item The collective organization shall have the right to pass simple resolutions to oversee resources belonging to the collective organization. This is a delegable right. Resources shall include the collective organization's:  
  \begin{enumerate}
    \item Funding and budgets.  
    \item Name and branding.  
    \item Public communications channels.  
    \item Space access.  
    \item Room reservations.  
    \item Property.  
    \item Secured documents.  
    \begin{enumerate}
      \item Retrieval or release of secured documents, particularly those produced within committees, officer boards, voting boards, or other private contexts, must further compelling interests of the collective organization.  
    \end{enumerate}
  \end{enumerate}
  \item The collective organization shall have the right to pass simple resolutions to regulate internal communications. This is a delegable right.  
  \begin{enumerate}
    \item Restrictions to internal communications must be narrowly tailored to further compelling interests of the collective organization.
  \end{enumerate}
\end{enumerate}

\subtab{Collective organization meetings}

\begin{enumerate}
  \item A collective organization meeting, or COM, shall be defined as a meeting which all members may attend as a right of membership. The primary purpose of COMs is to make decisions requiring the consensus of the organization at large. This definition is duplicated in section 2 ("definitions").  
  \item During fall and spring semesters, a COM must be scheduled no further than 14 days from the previous COM or from the start of the semester. One or more COMs must be scheduled within the last 14 days of the semester. The recommended spacing of COMs is every 7 days.  
  \item COMs should be scheduled outside of fall and spring semesters as required.
  \item The time and location of COMs must be announced with a minimum notice of 24 hours, barring unusual circumstances, and in a manner which all members can be reasonably expected to see.  
  \item COMs require quorum, which shall be defined as presence of all of the following:  
  \begin{enumerate}
    \item 1 or more clerks.  
    \item 2 or more officers, including clerks.  
    \item 7 or more members, including officers.  
  \end{enumerate}
  \item If a COM is scheduled and quorum is not reached for a minimum of 40 consecutive minutes or for as long as is necessary to complete all required business, whichever is shorter, another COM must be scheduled as soon as reasonably possible.  
  \item During a COM, some business items are required to occur. If one or more required items do not occur due to time constraints or exceptional circumstances, another COM must be scheduled as soon as reasonably possible as a continuation of the first, if reasonably necessary. The required items are:  
  \begin{enumerate}
    \item A statement by the officers describing ground rules and norms for the meeting, before any items which follow.  
    \item Officer reports, during which any information relevant to the collective organization relating to an officer position must be related by an officer holding that position, or by another individual if no such officer is present.  
    \begin{enumerate}
      \item Any changes to the officer board's composition during the period since the last COM must be reported during officer reports, whether by the voting board of a general or special election or by the officers.  
    \end{enumerate}
    \item Committee reports, during which any information relevant to the collective organization relating to committees' activities must be related by a chair of the committee, or by another individual if no such committee chair is present.  
    \begin{enumerate}
      \item Any expiration, dissolution, or changes to name, purpose, chairing, at-large membership requirements, or lifespan of any committee during the period since the last COM must be reported during committee reports. If the committee has regularly scheduled meetings open to members of the collective organization, the times and locations of upcoming meetings should be shared.  
    \end{enumerate}
    \item A period for other business to be conducted, during which any suitable rules of order may be used. This period must last for a minimum of 20 minutes, or until all business is completed, whichever occurs first.  
  \end{enumerate}
  \item Recommended, but not required, business items include the following:  
  \begin{enumerate}
    \item Introductions by all individuals present, which should include a name, pronouns, officer position(s), and any committees the individual is a chair or member of.  
    \item A time period during which all individuals present may list any upcoming actions or events they are aware of which may be of interest to the collective organization.  
  \end{enumerate}
  \item Business items, which include proposed resolutions, may be submitted in advance of a COM. All submitted business items must be brought up during the COM if reasonably possible, and the clerks should take reasonable steps to distribute time equitably among items to be considered.  
  \begin{enumerate}
    \item Intentional delaying or filibustering of COMs by individuals present shall be grounds for removal from the meeting, or for suspension of members if reasonable.
  \end{enumerate}
\end{enumerate}

\subtab{Resolutions}

\subsubtab{Proposing resolutions}

\begin{enumerate}
  \item Resolutions require one of the following forms of consensus to be considered formally proposed:
  \begin{enumerate}
    \item Two voting members may propose and sign off on the full contents of the amendment.
    \item Three voting and/or non-voting members may propose and sign off on the full contents of the amendment.
  \end{enumerate}
  \item The above forms of consensus apply both within COM and outside of it.
  \item Rules of procedure adopted by the collective organization may designate supplementary processes to propose resolutions during COM.
\end{enumerate}

\subsubtab{Passing resolutions}

Resolutions may be passed only by the following procedures:

\begin{enumerate}
  \item A simple resolution may be passed through a vote to pass by acclamation, which shall consist, at minimum, of a period of at least three seconds during which any voting member may voice an objection. If no objection is voiced, the resolution passes.  
  \item A ballot item in a formal vote.  
  \begin{enumerate}
    \item  The minimum notice period of the ballot item shall be 12 hours unless otherwise specified.  
    \item The ballot item shall include precisely three choices: "yes"/"pass", "no"/"fail", and "abstain" (or any other unambiguous naming scheme).  
    \item Vote counting shall discount all abstain votes, and the item shall pass if the required consensus is reached among "yes" and "no" votes.  
    \begin{enumerate}
      \item A simple resolution shall require majority consensus.  
      \item A supermajority resolution shall require supermajority consensus.
    \end{enumerate}
  \end{enumerate}
\end{enumerate}

\tab{Committees}

\subtab{Structure of committees}

A committee shall be defined as a group which coordinates activities relating to a specific, narrow, and time-sensitive purpose (this definition is duplicated in section 2, "definitions"). Committees shall have a limited lifespan. Committees shall formally consist of a small number of chairs initially chosen by the collective organization, who may oversee a broader committee membership. Committees must meet regularly.

\subtab{Committee formation}

Committees shall be created by the following process:

\begin{enumerate}
  \item A resolution to propose the creation of the committee must pass. This resolution must include:  
  \begin{enumerate}
    \item Its name and purpose.  
    \item Its initial number of chairs, which is to be used if the committee's chairs must be decided by ranked choice formal vote.  
    \begin{enumerate}
      \item The recommended number of chairs is 3 to 7. There must be no fewer than 2 chairs. There must be no more than 8 chairs, unless the resolution includes a clause permitting a larger maximum number of chairs.  
    \end{enumerate}
    \item Optionally, a clause increasing the maximum number of chairs, which may be raised to no more than 16.  
    \begin{enumerate}
      \item The initial size of the committee may be smaller than this bound, such as if there is an expectation that the committee may add chairs later on.  
    \end{enumerate}
    \item Its lifespan, which may be (i) a finite duration, (ii) a specific expiry date, or (iii) a specific event or occurrence after which the committee will expire. The lifespan is recommended to be approximately 28 days, and may not exceed 42 days, unless the resolution includes a clause permitting a longer maximum duration.  
    \item Optionally, a clause increasing the maximum duration of the committee's lifespan, which shall be no more than 365 days. Such a clause requires the resolution to pass by a majority vote.  
    \begin{enumerate}
      \item The initial lifespan of the committee may be less than this bound, such as if there is an expectation that the committee's lifespan may need to be extended.  
    \end{enumerate}
    \item Optionally, a clause delaying the formation of the committee. Such a clause must specify (i) a finite delay period, (ii) a specific date on which the committee shall be created, or (iii) a specific event or occurrence after which the committee shall be created.  
    \begin{enumerate}
      \item Unless the creation of the committee is tied to the expiry of another committee, if a committee's creation is delayed by further than 42 days, it shall no longer be created.  
      \item Time which passes during the delay period shall not count toward limits on the lifespan of a committee.  
    \end{enumerate}
  \end{enumerate}
  \item The initial committee chairs shall be determined by one of the following processes:  
  \begin{enumerate}
    \item A complete slate of committee chairs may be passed through a special resolution which may be passed by acclamation or through a supermajority formal ballot item. The size of the slate may differ from the fixed size chosen for ranked choice selection of chairs.  
    \item A ranked choice formal ballot item with STV vote counting may be used to elect committee chairs. Candidates shall be any members of the organization who self-nominate and are seconded by another member, or who are nominated by another member and accept nomination.  
  \end{enumerate}
  \item Any combination of formal ballot items required for the above process may be placed on the ballot of a single formal vote. If the results of the vote are contradictory, processes defined higher in this section shall take precedence.
\end{enumerate}

\subtab{Changes to existing committees}

\begin{enumerate}
  \item A committee may use any informal procedure to oversee its at-large membership.
  \begin{enumerate}
    \item Committees must abide by the CMU Statement of Assurance.
    \item Hazing must not be used as a condition of membership to any committee.
    \item Committee membership must not require dues or any form of financial commitment.
    \item Committee members must be free to voluntarily withdraw membership at any time.
    \item Any member of the organization must be able to join a committee, and to do so within a reasonable time period, unless the committee is not accepting members or has a compelling reason to deny membership on an individual basis. The officer board may override committee decisions on at-large membership in cases where the committee cannot provide justification for rejecting membership or removing a member.
  \end{enumerate}
  \item A committee may use any informal procedure internally to add one or more chairs, subject to its maximum number of chairs.  
  \begin{enumerate}
    \item Should there be internal disagreement on the addition of a chair to a committee, any chair of the committee may present a written petition for a vote, which shall require a written majority of committee chairs to support the addition of the chair, subject to enforcement by the officers of the organization.  
  \end{enumerate}
  \item Chairs may leave a committee voluntarily through written notice.  
  \begin{enumerate}
    \item The minimum number of committee chairs shall be 2. If at any point this is violated, the committee is dissolved immediately.  
  \end{enumerate}
  \item A committee may remove a chair through the written consent of a supermajority of other chairs of the committee, with the written consent of a supermajority of officers.  
  \begin{enumerate}
    \item The minimum number of committee chairs shall be 2. If at any point this is violated, the committee is dissolved immediately.  
  \end{enumerate}
  \item A committee may propose a change to its name, purpose, or both, which shall be approved by a simple resolution of the collective organization.  
  \item A committee may change its lifespan, providing that it does not exceed the committee's maximum lifespan, which is 42 days unless increased by the collective organization.  
  \item A simple resolution through the collective organization may alter any of the following:  
  \begin{enumerate}
    \item The maximum number of chairs a committee may have, up to no more than 16.  
    \item The maximum lifespan of a committee, up to no more than 365 days.  
  \end{enumerate}
  \item A supermajority resolution through the collective organization may alter any of the following:  
  \begin{enumerate}
    \item The chairs of a committee.  
    \item The lifespan of a committee.  
  \end{enumerate}
  \item A committee may opt to dissolve itself.  
  \item A supermajority resolution may dissolve a committee.
\end{enumerate}

\subtab{Committee meetings}

\begin{enumerate}
  \item A committee meeting shall be defined as a meeting at which a committee's chairs are present and which all at-large members of the committee may attend as a right of membership (this definition is duplicated in section 2, "definitions").
  \item Committees must schedule their first committee meeting within 7 days of the committee's formation, if reasonably possible.
  \item Committees must schedule committee meetings no further than 21 days from any previous committee meeting if reasonably possible.
  \item Multiple committees may conduct joint committee meetings, if attendees will have a reasonable opportunity to engage fully with each committee present.
  \item Committee meetings must be announced such that all chairs and at-large members can be reasonably expected to be aware of the time and location of the meeting.
  \item Committees are encouraged to make meetings open to the collective organization's membership, and committees which choose to do so should announce meetings in such a manner that it can be reasonably expected that any member of (i) one or more relevant interest groups or (ii) any member of the collective organization will be aware of the time and location of the meeting.
\end{enumerate}

\tab{Interest groups}

\subtab{Interest groups}

\begin{enumerate}
  \item An interest group, or IG, shall be defined as a group of individuals and one or more committees with a shared focus, with the purpose of sharing information and updates relating to that focus, rather than carrying out actions, which is the duty of committees. This definition is duplicated in section 2 ("definitions").  
  \item An interest group must include one committee which is assigned the role of "administering committee", which is responsible for:  
  \begin{enumerate}
    \item Overseeing the composition of the interest group.  
    \item Communicating to the interest group any relevant public committee meetings or announcements.  
  \end{enumerate}
  \item An interest group may be created through a resolution. Such a resolution must designate a committee, new or existing, to serve as the administering committee.  
  \begin{enumerate}
    \item If a new committee is part of a proposed interest group, the resolution must include all details applicable to a resolution to create a committee.  
    \item If an existing committee is part of a proposed interest group, a chair of the committee must affirm the committee's interest in, and capability to, administer the interest group.  
  \end{enumerate}
  \item An interest group may exist indefinitely, should it continue to have an administering committee. If an interest group has no administering committee for more than 14 days, it is immediately dissolved.  
  \item An interest group may include committees other than its administering committee, which may be added or removed by the administering committee. A committee must consent to its inclusion in an interest group.  
  \item With another committee's consent, the administering committee may (i) transfer its status to that committee, or (ii) designate that committee as a numbered successor should the current administering committee expire or otherwise dissolve.  
  \item Interest groups may also include individuals, who must be members of the organization. The administering committee has a duty to ensure all interest group members are organization members, and that all organization members may join the interest group if desired.  
  \item The administering committee of an interest group may be changed by a supermajority resolution.  
  \item The administering committee of an interest group may opt to dissolve it.  
  \item An interest group may be dissolved by a supermajority resolution.
\end{enumerate}

\subtab{Special interest groups}

\begin{enumerate}
  \item A special interest group, or SIG, shall be defined as an interest group with a limited lifespan. This definition is duplicated in section 2 ("definitions").  
  \item The resolution which proposes a special interest group must include (i) a finite duration of the special interest group, (ii) a specific date on which the special interest group will expire, or (iii) a specific event or occurrence after which the special interest group will expire.  
  \item The lifespan of a special interest group may be changed through a resolution. The administering committee may delay the expiry of the special interest group until the next collective organization meeting, and if a resolution to extend the lifespan of the special interest group goes to a formal vote, the SIG's expiry is delayed until the formal vote occurs and votes are counted.
\end{enumerate}

\tab{Voting}

\subtab{Written votes}

\begin{enumerate}
  \item A written vote shall be defined as a procedure to decide a question among a group of fixed size, such as a committee or officer board. This definition is duplicated in section 2 ("definitions").  
  \item The procedure to conduct a written vote shall be as follows:  
  \begin{enumerate}
    \item Whichever individual called for the written vote must specify a reasonable time constraint under which the vote is to occur.  
    \item Within this time, voters should submit their vote in a manner which may be documented. The procedure may be anonymized if it occurs through in-person paper ballots, but otherwise must have voters' identities attached.  
    \item Votes must be counted and documented following the completion of the voting period or the submission of votes by all eligible voters.  
  \end{enumerate}
  \item The officer board reserves the right to maintain the outward anonymity of individual officers' votes.
\end{enumerate}

\subtab{Formal votes}

\subsubtab{Purpose of formal votes}

A formal vote shall be defined as an anonymous process to decide one or more questions through an in-person election. Formal votes shall occur during COMs. Most COMs will occur to resolve proposed resolutions from the previous COM, or for general or special elections of officers. This definition is duplicated in section 2 ("definitions").

\subsubtab{Procedure}

A formal vote shall occur by the following procedure:

\begin{enumerate}
  \item Formal votes shall be overseen by a voting board. The voting board should be decided once for each vote which occurs. The membership may be decided by any informal process, but must consist of voting members. If at any point the voting board does not have two or more members, it must be dissolved and rebuilt in a timely manner.  
  \begin{enumerate}
    \item For any formal votes containing ballot items which will decide officer positions, no candidate for such a position shall stand on the voting board, unless no reasonable alternative is possible due to membership or scheduling limitations.  
  \end{enumerate}
  \item Each ballot item on a formal vote's ballot shall have an associated minimum notice period. The voting board must ensure an announcement of the ballot item's contents is made prior to this notice period in a manner which all voting members can reasonably be expected to see. As soon as a resolution is formally proposed, assuming it is not passed by acclamation, a formal vote must be scheduled at the soonest or second-soonest possible COM which maintains sufficient notice.  
  \item Formal votes shall occur during COMs. During a formal vote, a ballot box must be placed in plain view of all attendees of the COM for the full duration of a voting period which is sufficiently long so as to allow all voting members present to finish voting without undue hurry. The voting period must start no later than 40 minutes past the scheduled start of the COM. The ballot box must be watched at all times by one or more members of the voting board.  
  \item The voting board must record the name, Andrew ID, and/or any other information necessary for unambiguous identification of all voters on a physical or digital document of some form, which must be visible to any watchers throughout the process of voting and counting votes.  
  \item The voting board must take reasonable steps to prevent any member from voting two or more times and to prevent non-members and non-voting members from voting.  
  \item In order to cast a vote, a voter may do one of the following:  
  \begin{enumerate}
    \item Write or otherwise indicate their choices on a card or piece of paper, in a manner chosen by the voting board. The voting board must make such cards or standardized pieces of paper available to voters, as well as writing utensils or some acceptable method of marking ballots.  
    \item Submit an absentee ballot to the voting board prior to or during the voting period, in a manner which permits the voting board to unambiguously identify the voter. Anonymity shall not be possible with this process, and all members of the voting board or ballot watchers may inspect the means by which absentee ballots are validated and counted.  
  \end{enumerate}
  \item With the exception of absentee voting, voting must be anonymous. The voting process must not be implemented in such a manner as to make it reasonably possible to identify the voter who submitted a particular physical ballot.  
  \item Immediately following the conclusion of the voting period, the voting board must open the ballot box and record all ballots' contents in a manner which can be stored permanently, whether physically or digitally. The voting board should take reasonable steps to mix or shuffle ballots such that their order or placement cannot be used to associate ballots with voters. Any member shall have the right to watch ballots be counted.  
  \begin{enumerate}
    \item Any illegible or ambiguous responses to a ballot item must not be counted.  
    \item The recorded contents of ballot items must be limited to the minimal information required for counting votes. No photographs of ballots or other information which may aid in identifying voters may be included.
  \end{enumerate}
  \item The voting board must destroy all ballots after counting. There may not be any period of time after the voting period and before the ballots are destroyed where one or more members of the voting board are not present. Any member shall have the right to watch the ballot box during this period.  
  \item Members of the voting board may vote. They must do so during the voting period and in the presence of one or more other members of the voting board. They must record their name in the same manner as all other voters.
\end{enumerate}

\tab{Amendments}

Amending these bylaws shall occur by the following process:

\begin{enumerate}
  \item A supermajority resolution must be passed through the collective organization. The ballot item must have a minimum notice period of 6 days.  
  \begin{enumerate}
    \item The ballot item must contain the precise text of the amendment as written in the resolution proposing the amendment. % Note: This is shitty wording I didn't notice while typing up this amendment in diff form. Fix this in a future amendment :-)
  \end{enumerate}
  \item Following the counting of votes for the ratification of an amendment, should it pass, the amended bylaws shall take effect immediately. An updated copy of the bylaws must be made available within 7 days.
\end{enumerate}

\tab{Dissolution}

\subtab{Voluntary dissolution}

Voluntary dissolution of the organization shall occur by the following process:

\begin{enumerate}
  \item A simple resolution to propose the dissolution of the organization must be passed through the collective organization.  
  \item Following this resolution, a formal vote must be scheduled to contain a supermajority ballot item to ratify the dissolution. The ballot item must have a minimum notice period of 6 days.  
  \item Following the counting of votes for the dissolution of the organization, should it pass, the organization shall become "undergoing dissolution" as defined in this section.
\end{enumerate}

\subtab{Dissolution process}

The status of "undergoing dissolution," which shall apply following voluntary dissolution or involuntary dissolution by an outside authority, shall impose the following requirements:

\begin{enumerate}
  \item The organization, its officers, and its committees must cease any activities conducted under the name or branding of the organization.  
  \item No property of the organization nor any of its monetary assets shall be distributed to any members of the organization. After payment of the debts of the organization, its property and monetary assets must be distributed as required by outside authorities and with a preference toward the continued furtherance of the mission and purposes of the organization. Resolutions passed by the collective organization prior to its dissolution may influence this process and may designate an individual or group as having the right to conduct it.
\end{enumerate}

\tab{History}

\begin{enumerate}
  \item Revision 1.0: Ratified September 29, 2025.  
  \item Revision 1.1: Amended via formal vote November 20th, 2025.  
  \begin{enumerate}
    \item Consists of amendments 1, 2, and 3, all proposed on November 9th, 2025.
  \end{enumerate}
  \item Revision 1.2: Amended via formal vote February 25th, 2025. Consists of:
  \begin{enumerate}
    \item Amendments 4, 5, 6, 7, and 8, all proposed on February 5th, 2025, and
    \item Amendments 9, 10, 11, and 12, all proposed on February 12th, 2025.
  \end{enumerate}
\end{enumerate}

\end{document}